% appendix-lower: proofs for the lower bounds of the complexity section
\section{Proofs of the Lower Bounds}
\label{app:lower}

All trees in this appendix are covariant, and $n\ge2$. We use the notation
of \cref{sec:complexity-lower}. For a word $v$, $\pr(v)$ denotes the set of its
prefixes, including $v$; $\pr(v)$ is finite and prefix closed.

\subsection{Top-Free Trees}

\begin{proof}[Proof of \cref{lem:cx-topfree}]
We first note that $\widehat t(w)=f$ iff $w\in\Pos(t)$, for every tree $t$ at
positive arity. If $w\in\operatorname{dom}(t)$, then every proper prefix of $w$
is labeled $f$ by \eqref{eq:g2}, so the walk along $w$ reaches its end and
$\widehat t(w)=t(w)$. Otherwise, the longest prefix of $w$ in
$\operatorname{dom}(t)$ is a leaf by \eqref{eq:g2}, and a leaf is labeled $\bot$ or
$\top$.

(1) For a top-free tree, leaves are labeled $\bot$, so by the remark
$\widehat t(w)=\bot$ for $w\notin\Pos(t)$. For top-free $t,u$ both leaf extensions
take values in $\{\bot,f\}$, so by \eqref{eq:g11} $t\le u$ holds iff
$\widehat t(w)=f$ implies $\widehat u(w)=f$ for every $w$, that is, iff
$\Pos(t)\subseteq\Pos(u)$.

(2) The set $\{\eps\}\cup P\cup P\Dir$ contains the root and is prefix closed,
because $P$ is. For a word $w$ in it and $i\in\Dir$, the word $wi$ belongs to it
iff $w\in P$: if $wi\in P$, then $w\in P$ by prefix closure. So labeling $P$ by
$f$ and the other words by $\bot$ satisfies \eqref{eq:g2}, and $\Pos(T_P)=P$.
A top-free tree $t$ is determined by $\Pos(t)$, since by \eqref{eq:g2} its
domain is $\{\eps\}\cup\Pos(t)\cup\Pos(t)\Dir$; hence $T_P$ is unique. The
domain is finite iff $P$ is, since $\Dir$ is finite.

(3) Write $u=\sigma(g)$ and $s=f(u,u,\Bot,\ldots,\Bot)$, using $\sigma(b)=\Bot$.
Then $\widehat s(\eps)=f$, $\widehat s(iw)=\widehat u(w)$ for $i\in\{0,1\}$, and
$\widehat s(iw)=\bot$ for $2\le i<n$. Since $u\le s$, \eqref{eq:g11} gives
$\widehat u(w)\le\widehat s(w)$ for every $w$. By induction on $|w|$ we show that
$\widehat u(w)\ne\top$, and that $w\in\Pos(u)$ implies $w\in\{0,1\}^*$. For
$w=\eps$, $\widehat u(\eps)\le f$. For a word $iw$ with $i\ge2$,
$\widehat u(iw)\le\bot$, so $\widehat u(iw)=\bot$ and $iw\notin\Pos(u)$. For
$i\in\{0,1\}$, $\widehat u(iw)\le\widehat u(w)\ne\top$; if moreover $iw\in\Pos(u)$,
then $f=\widehat u(iw)\le\widehat u(w)$ forces $\widehat u(w)=f$, so $w\in\Pos(u)$
and, by induction, $iw\in\{0,1\}^*$. A node labeled $\top$ at $w$ would give
$\widehat u(w)=\top$, so $u$ is top free. For $t\le s$, the same comparison gives
$\widehat t(w)\le\widehat s(w)\ne\top$ for every $w$, so $t$ is top free; and if
$iw\in\Pos(t)$, then $f\le\widehat s(iw)$ forces $i\in\{0,1\}$ and
$\widehat u(w)=f$, so $iw\in\{0,1\}\Pos(u)\subseteq\{0,1\}^*$.

(4) Following a word $w$ in $T_P^m$ gives
\[
 \widehat{T_P^m}(w)=
 \begin{cases}
 \top,&\len w\ge m\text{ and }\wpre wm\in P,\\
 f,&\len w<m\text{ and }w\in P,\\
 \bot,&\text{otherwise}.
 \end{cases}
\]
Let $s=f(t_0,\ldots,t_{n-1})$. By \eqref{eq:g11} it suffices to show
$\widehat s(w)\le\widehat{T_P^m}(w)$ for every $w$. At $\eps$ both sides are
$f$, since $\eps\in P$ and $m\ge1$. At $iv$ with $i\ne a$ the left side is
$\bot$. At $av$ the left side is $\widehat{T_{P'}^m}(v)$. If it is $\top$, then
$\len v\ge m$ and $\wpre vm\in P'$, so $a\wpre vm\in P$, hence
$\wpre{(av)}m\in P$ by prefix closure, and the right side is $\top$. If it is
$f$, then $\len v<m$ and $av\in P$; the right side is $f$ if $\len{av}<m$, and
$\top$ if $\len{av}=m$. Otherwise the left side is $\bot$.
\end{proof}

\subsection{The Counter Automaton}

\begin{proof}[Proof of \cref{lem:cx-counter}]
\emph{A token automaton.} Let $\mathcal N_\beta$ be the nondeterministic
automaton over the token alphabet $\mathbb T$ with the following states and transitions; a transition that
is not listed does not exist. Bits are the letters $0,1$.
\begin{itemize}[leftmargin=*]
\item $\mathsf E$ loops on every letter.
\item $\mathsf F_0,\ldots,\mathsf F_\beta$ check the format: $\mathsf F_i$ goes on a
  bit to $\mathsf F_{i+1}$ if $i<\beta$ and to $\mathsf E$ if $i=\beta$; on $\#$ to
  $\mathsf F_0$ if $i=\beta$ and to $\mathsf E$ otherwise; and on $\$$ to $\mathsf E$
  if $i<\beta$.
\item $\mathsf I$ checks that the first block is zero: it goes on $0$ to
  $\mathsf I$ and on $1$ to $\mathsf E$.
\item $\mathsf C_\gamma$, for a carry $\gamma\in\{0,1\}$, tracks the carry of an
  increment and spawns checkers: on a bit $a$ it goes both to
  $\mathsf C_{\gamma\wedge a}$ and to $\mathsf V_{a\oplus\gamma,\beta}$, and on $\#$ to
  $\mathsf C_1$.
\item $\mathsf V_{a,j}$, for $a\in\{0,1\}$ and $0\le j\le\beta$, waits for the
  corresponding bit of the next block: for $j>0$ it goes on every letter other
  than $\$$ to $\mathsf V_{a,j-1}$, and $\mathsf V_{a,0}$ goes on the bit $1-a$ to
  $\mathsf E$.
\item $\mathsf O_o$, for $o\in\{0,1\}$, records whether the current block is all
  ones so far: on a bit $a$ it goes to $\mathsf O_{o\wedge a}$; on $\#$ to
  $\mathsf E$ if $o=1$ and to $\mathsf O_1$ if $o=0$; and $\mathsf O_0$ goes on $\$$
  to $\mathsf E$.
\end{itemize}
The initial states are $\mathsf F_0$, $\mathsf I$, $\mathsf C_1$, and $\mathsf O_1$, and
there are $3\beta+9$ states, all final. The following table summarizes the
roles of the state families.
\begin{center}
\small
\begin{tabular}{@{}lp{.36\linewidth}p{.36\linewidth}@{}}
\toprule
States & Error detected & Behavior on a prefix of $\mathsf H_\beta$\\
\midrule
$\mathsf E$ & none; records an error found earlier & never reached\\
$\mathsf F_0,\ldots,\mathsf F_\beta$ & a block with other than $\beta$ bits & counts the bits of each block; dies at $\$$\\
$\mathsf I$ & a bit $1$ in the first block & reads the first block; dies at the first $\#$\\
$\mathsf C_0,\mathsf C_1$ & none; spawns the checkers & tracks the carry of the increment; dies at $\$$\\
$\mathsf V_{a,j}$ & a wrong bit $\beta+1$ tokens after the bit that spawned it & dies on reading the predicted bit, or at $\$$\\
$\mathsf O_0,\mathsf O_1$ & $\#$ after the all-ones block, or $\$$ after another block & dies at $\$$\\
\bottomrule
\end{tabular}
\end{center} We show that the words with a run of
$\mathcal N_\beta$ are exactly the words of $\mathbb T^*\setminus\mathsf H_\beta\mathbb T^*$. Write
$\mathsf H=\mathsf H_\beta$; the bit of $\cxbin j$ at offset $p<\beta$ is at
position $j(\beta+1)+p$ of $\mathsf H$.

No branch reaches $\mathsf E$ while reading a prefix of $\mathsf H$. The format
is correct. The first block is $0^\beta$. Along $\mathsf H$ the carry branch is
deterministic: after the bits of $\cxbin j$ at offsets below $p$, it is in
$\mathsf C_\gamma$ with $\gamma$ the carry into offset $p$ when $1$ is added to
$j$. On the bit $a$ at offset $p$ it therefore spawns
$\mathsf V_{a\oplus\gamma,\beta}$. For $j<2^\beta-1$, the bit $a\oplus\gamma$ is
the bit of $\cxbin{j+1}$ at offset $p$. That bit lies $\beta+1$ positions later,
so the checker reads it in state $\mathsf V_{a\oplus\gamma,0}$ and has no
transition on it. A $\#$ follows only blocks that are not all ones, and $\$$
follows the all-ones block $\cxbin{2^\beta-1}$. After $\mathsf H$, no state remains:
$\mathsf F_\beta$, the states $\mathsf C_\gamma$, and $\mathsf O_1$ have no
transition on $\$$, $\mathsf I$ has none on the first $\#$, the checkers of
earlier blocks have died as just shown, and a checker spawned at offset $p$ of
the last block reads $\$$ in state $\mathsf V_{a,p+1}$, which has no transition
on $\$$. So no word of $\mathsf H\mathbb T^*$ has a run.

Every proper prefix of $\mathsf H$ has a run through the format states. Let
$u\notin\mathsf H\mathbb T^*$ be no prefix of $\mathsf H$, and let $i<\len{\mathsf H}$ be
the first position with $u[i]\ne\mathsf H[i]$. Some branch that is active after
$\wpre ui=\wpre{\mathsf H}i$ reaches $\mathsf E$ on $u[i]$:
\begin{itemize}[leftmargin=*]
\item if $\mathsf H[i]$ is the bit at offset $p$ of $\cxbin j$ and $u[i]$ is
  $\#$ or $\$$, the format branch is in $\mathsf F_p$ with $p<\beta$;
\item if $u[i]$ is the other bit, then for $j=0$ the branch $\mathsf I$ reads
  $1$, and for $j\ge1$ the checker spawned at offset $p$ of $\cxbin{j-1}$ is in
  $\mathsf V_{\mathsf H[i],0}$;
\item if $\mathsf H[i]=\#$, a bit leads $\mathsf F_\beta$ to $\mathsf E$, and $\$$
  leads $\mathsf O_0$ to $\mathsf E$, since the block before is not all ones;
\item if $\mathsf H[i]=\$$, a bit leads $\mathsf F_\beta$ to $\mathsf E$, and $\#$
  leads $\mathsf O_1$ to $\mathsf E$.
\end{itemize}
Since $\mathsf E$ loops, $u$ has a run, which proves the claim.

\emph{Binary words.} Let $\mathcal B_\beta$ have the states $q$, $(q,0)$, and
$(q,1)$ for every state $q$ of $\mathcal N_\beta$, all final, the initial states
of $\mathcal N_\beta$, a transition on the bit $a$ from $q$ to $(q,a)$, and a
transition on the bit $a'$ from $(q,a)$ to every successor of $q$ in
$\mathcal N_\beta$ on the token encoded by $aa'$. It has $9\beta+27$
states. By induction on the number of encoded letters, the states of the form
$q$ reachable on the token encoding of a word $z\in\mathbb T^*$ are the states of
$\mathcal N_\beta$ reachable on $z$, and after one further bit $a$ the reachable
states are the $(q,a)$ for these $q$; this set is nonempty iff the former is.
So a binary word $u$ has a run iff the word $z\in\mathbb T^*$ encoded by the longest
prefix $u'$ of $u$ of even length has a run of $\mathcal N_\beta$, that is, iff
$z\notin\mathsf H\mathbb T^*$. As $\mathsf W_\beta$ has even length and the encoding uses
blocks of two bits, $z\in\mathsf H\mathbb T^*$ iff $u'\in\mathsf W_\beta\{0,1\}^*$ iff
$u\in\mathsf W_\beta\{0,1\}^*$. The construction is explicit and takes time
polynomial in $\beta$.
\end{proof}

\subsection{The Embedding of Binary Trees}

The remark on padding in \cref{sec:complexity-lower} compares systems over
a binary constructor with their padded versions by an embedding of binary
trees; no other result depends on it. Let
$n\ge3$, let $f$ be covariant of arity $n$, and let $f_2$ be binary and
covariant. For a binary tree $t$, let $\cxlift(t)$ be the tree over $\Dir$
with domain $\operatorname{dom}(t)\cup\{wi:w\in\Pos(t),\ 2\le i<n\}$, the
labels of $t$ on $\operatorname{dom}(t)$ with $f_2$ renamed to $f$, and $\Bot$
elsewhere. For a tree $t'$ over $\Dir$, let $\cxproj(t')$ be its restriction to
$\operatorname{dom}(t')\cap\{0,1\}^*$ with $f$ renamed to $f_2$. For a binary
tree, $\Pos$ collects the positions labeled $f_2$, and leaf extensions of
trees over the two signatures are compared with $f_2$ identified with $f$.
For a flat system $\chi$ over $f_2$ and a variable $o$, either with the
literal $o=\Bot$ in $\chi$ or fresh, let $\chi^{(n)}$ replace every
$f_2(a_0,a_1)$ by $f(a_0,a_1,\bar o)$ and contain the literal $o=\Bot$.

\begin{proposition}[Embedding binary trees]\label{prop:cx-embedding}
Let $n\ge3$, and let $f$, $f_2$, $\cxlift$, $\cxproj$, and $\chi^{(n)}$ be as
above.
\begin{enumerate}[leftmargin=*]
\item $\cxproj\circ\cxlift$ is the identity, $\cxlift$ and $\cxproj$ are
  monotone and preserve finite and regular trees, and
  $\widehat{\cxlift(t)}(w)=\widehat t(w)$ and
  $\widehat{\cxproj(t')}(w)=\widehat{t'}(w)$ for every $w\in\{0,1\}^*$.
\item In each domain, $\cxlift\circ\sigma$, extended by $o\mapsto\Bot$, solves
  $\chi^{(n)}$ whenever $\sigma$ solves $\chi$, and $\cxproj\circ\sigma'$ solves
  $\chi$ whenever $\sigma'$ solves $\chi^{(n)}$. Hence $\chi$ and $\chi^{(n)}$
  are satisfiable in the same domains, and for every query $x\le y$ a word
  in $\{0,1\}^*$ is left (right) unsafe for $\chi$ iff it is left (right)
  unsafe for $\chi^{(n)}$.
\item If $\chi$ contains $g\le f_2(g,g)$ and $x\le f_2(g,g)$, then no word is
  left unsafe for $(\chi^{(n)},x,y)$, and every right-unsafe word lies in
  $\{0,1\}^*$.
\end{enumerate}
\end{proposition}

\begin{proof}[Proof of \cref{prop:cx-embedding}]
As above, $f_2$ is identified with $f$ whenever labels or leaf
extensions of trees over the two signatures are compared.

(1) For a binary tree $t$, the domain of $\cxlift(t)$ is prefix closed and
contains $\eps$. It satisfies \eqref{eq:g2}: for $w\in\operatorname{dom}(t)$, a
child $wi$ with $i\in\{0,1\}$ belongs to it iff $t(w)=f_2$; a child $wi$ with
$2\le i<n$ belongs to it iff $w\in\Pos(t)$; and the added nodes are labeled
$\bot$ and have no children. Similarly, $\operatorname{dom}(t')\cap\{0,1\}^*$
satisfies \eqref{eq:g2} for $f_2$, so $\cxproj(t')$ is a binary tree. As
$\operatorname{dom}(\cxlift(t))\cap\{0,1\}^*=\operatorname{dom}(t)$,
$\cxproj\circ\cxlift$ is the identity. For $w\in\{0,1\}^*$, the walk along $w$
visits the same nodes, with the same labels up to renaming, in $t$ and
$\cxlift(t)$, and in $t'$ and $\cxproj(t')$, so the leaf extensions agree at
$w$. For monotonicity use \eqref{eq:g3}: the common positions of $\cxlift(t)$ and $\cxlift(u)$ in
$\{0,1\}^*$ are common positions of $t$ and $u$ with the same labels, and the
others end in a letter $i\ge2$ and are labeled $\bot$ in both; the common
positions of $\cxproj(t')$ and $\cxproj(u')$ are common positions of $t'$ and
$u'$ with the same labels. Both maps preserve finite domains. Every subtree of
$\cxlift(t)$ is the image of a subtree of $t$ or is $\Bot$, and every subtree of
$\cxproj(t')$ is the image of a subtree of $t'$, so both maps preserve
regularity.

(2) Both maps fix $\Bot$ and $\Top$, and
\[
 \cxlift(f_2(t_0,t_1))=f(\cxlift t_0,\cxlift t_1,\Bot,\ldots,\Bot),\qquad
 \cxproj(f(t_0,\ldots,t_{n-1}))=f_2(\cxproj t_0,\cxproj t_1).
\]
Let $\sigma$ solve $\chi$. Applying the monotone map $\cxlift$ to a literal
$f_2(\sigma a_0,\sigma a_1)\le\sigma z$ gives
\[
 f(\cxlift\sigma a_0,\cxlift\sigma a_1,\Bot,\ldots,\Bot)\le\cxlift\sigma z,
\]
which is the translated literal under $\cxlift\circ\sigma$ with $o\mapsto\Bot$; if $o$ is
a variable of $\chi$, then $\sigma(o)=\Bot$ and $\cxlift(\Bot)=\Bot$. Upper and
constant literals are treated in the same way, and $o=\Bot$ holds. Conversely,
applying $\cxproj$ to a translated literal under a solution $\sigma'$ of
$\chi^{(n)}$ deletes the arguments $2,\ldots,n-1$ and gives the original literal
under $\cxproj\circ\sigma'$. By~(1), both maps preserve each domain. For a binary
word $w$, the leaf extensions at $w$ are preserved by~(1), so a solution that
witnesses $w$ as left or right unsafe for one system is mapped to a solution
that witnesses the same for the other.

(3) Let $\sigma'$ solve $\chi^{(n)}$. The literals $o=\Bot$,
$g\le f(g,g,\bar o)$, and $x\le f(g,g,\bar o)$ and \cref{lem:cx-topfree}(3), with
$o$ in the role of $b$, show that $\sigma'(x)$ is top free with
$\Pos(\sigma'(x))\subseteq\{0,1\}^*$. Hence $\widehat{\sigma'(x)}(w)\ne\top$ for
every $w$, and if $\widehat{\sigma'(x)}(w)\ne\bot$ then
$w\in\Pos(\sigma'(x))\subseteq\{0,1\}^*$ by \cref{lem:cx-topfree}(1).
\end{proof}

\subsection{Rejection Systems}

\begin{proof}[Proof of \cref{lem:cx-rejection}]
Let $\mathcal A=(Q,I,Q,\delta)$ be all-final, and for
$q\in Q\cup\{q_{\mathrm{in}}\}$ let $L_q$ be the set of words that label a
path from $q$. Each $L_q$ is prefix closed and contains $\eps$, and
$aL_{q'}\subseteq L_q$ for every transition $q\xrightarrow{a}q'$. Moreover
$L_{q_{\mathrm{in}}}=L(\mathcal A)$: a nonempty word has a path from
$q_{\mathrm{in}}$ iff it has one from some state of $I$, and $\eps$ belongs
to both sets because $I\ne\varnothing$. Recall that the systems are
covariant, so a word $w$ is left unsafe iff some solution $\sigma$ has
$\widehat{\sigma(x)}(w)=\top\ne\widehat{\sigma(y)}(w)$, and right unsafe iff
some solution has $\widehat{\sigma(y)}(w)=\bot\ne\widehat{\sigma(x)}(w)$.

(1) Put $\sigma(x)=\sigma(g)=\sigma(b)=\Bot$ and $\sigma(Y_q)=\Top$ for every
$q\in Q\cup\{q_{\mathrm{in}}\}$. The literal $b=\Bot$ holds, the two upper
literals have $\Bot$ on the left, and every lower literal has $\Top$ on the
right. These trees are finite.

\emph{Claim A.} Every solution $\sigma$ of $\Phi$ has $\sigma(x)$ top free
with $\Pos(\sigma(x))\subseteq\{0,1\}^*$, and
$\widehat{\sigma(Y_q)}(v)\ne\bot$ for every $q\in Q\cup\{q_{\mathrm{in}}\}$ and
$v\in L_q$. The first part is \cref{lem:cx-topfree}(3), applied to the
literals $b=\Bot$, $g\le f(g,g,\bar b)$, and $x\le f(g,g,\bar b)$. For the
second, use induction on $v$. For $v=\eps$, the literal
$f(b,b,\bar b)\le Y_q$ and \eqref{eq:g11} give
$\widehat{\sigma(Y_q)}(\eps)\ge f$. For $v=av'$, some transition
$q\xrightarrow{a}q'$ has $v'\in L_{q'}$. The left side $s$ of the
corresponding literal, say $f(Y_{q'},b,\bar b)\le Y_q$ for $a=0$, has
$\widehat s(0v')=\widehat{\sigma(Y_{q'})}(v')$, so \eqref{eq:g11} gives
$\widehat{\sigma(Y_q)}(0v')\ge\widehat{\sigma(Y_{q'})}(v')$, which is not
$\bot$ by induction; the case $a=1$ is symmetric.

(2) Let $\sigma$ be a solution. By Claim A and \cref{lem:cx-topfree}(1),
$\widehat{\sigma(x)}(w)\in\{\bot,f\}$ for every word $w$, so no word is left
unsafe. If $w$ is right unsafe, witnessed by $\sigma$, then
$\widehat{\sigma(x)}(w)\ne\bot$, so $w\in\Pos(\sigma(x))\subseteq\{0,1\}^*$ by
\cref{lem:cx-topfree}(1), and $\widehat{\sigma(y)}(w)=\bot$, so
$w\notin L_{q_{\mathrm{in}}}=L(\mathcal A)$ by Claim A.

(3) Let $w\in\{0,1\}^*\setminus L(\mathcal A)$ and $m=\len w+1$, and put
\[
 \sigma_w(b)=\Bot,\qquad \sigma_w(g)=T_{\{0,1\}^{<m}},\qquad
 \sigma_w(x)=T_{\pr(w)},\qquad \sigma_w(Y_q)=T_{L_q}^m\ \ (q\in Q\cup\{q_{\mathrm{in}}\}).
\]
These trees are finite. The literal $b=\Bot$ holds. The tree
$s=f(\sigma_w(g),\sigma_w(g),\Bot,\ldots,\Bot)$ is top free with
$\Pos(s)=\{\eps\}\cup\{0,1\}\{0,1\}^{<m}=\{0,1\}^{\le m}$. By
\cref{lem:cx-topfree}(1), $g\le f(g,g,\bar b)$ holds because
$\{0,1\}^{<m}\subseteq\{0,1\}^{\le m}$, and $x\le f(g,g,\bar b)$ holds
because $\pr(w)\subseteq\{0,1\}^{\le m}$. Each lower literal holds by
\cref{lem:cx-topfree}(4) with $P=L_q$: the literal $f(b,b,\bar b)\le Y_q$
is the case $P'=\varnothing$, where $T_\varnothing^m=\Bot$, and the literal
of a transition $q\xrightarrow{a}q'$ is the case $P'=L_{q'}$. Thus
$\sigma_w$ solves $\Phi$. Finally $\widehat{\sigma_w(x)}(w)=f$, since
$w\in\pr(w)=\Pos(\sigma_w(x))$, while $\len w<m$ and
$w\notin L_{q_{\mathrm{in}}}$ give $\widehat{\sigma_w(y)}(w)=\bot$ by the
formula in the proof of \cref{lem:cx-topfree}(4). So $w$ is right unsafe,
and $\sigma_w(x)\not\le\sigma_w(y)$ by \eqref{eq:g11}.

\emph{The equivalence.} By~(2) and~(3), the right-unsafe words are exactly
the words of $\{0,1\}^*\setminus L(\mathcal A)$. If $\mathcal A$ is
universal, no word is left or right unsafe, so $\Phi\models x\le y$ over
arbitrary trees by \eqref{eq:g13}. Regular and finite solutions are in
particular solutions over arbitrary trees, so the entailment holds in each
domain. If $\mathcal A$ rejects a word $w$, then $\sigma_w$ is a finite
solution, hence a regular one, with $\sigma_w(x)\not\le\sigma_w(y)$, and the
entailment fails in each domain.
\end{proof}
